\documentclass[aspectratio=169]{beamer}

\usetheme{Madrid}
\usecolortheme{default}

\usepackage{booktabs}
\usepackage{xcolor}
\usepackage{amsmath}
\usepackage{amssymb}
\usepackage[round]{natbib}

\definecolor{color}{rgb}{0.15294, 0.21176, 0.22475}
\definecolor{darkred}{rgb}{0.75, 0.0, 0.0}
\usecolortheme[named=color]{structure}

\title{Synchronic Semantic Variation Across Groups}
\subtitle{CSC2611 Final Presentation}
\author{Yimo Ning, Johnny Meng}
\date{27 March 2026}

\begin{document}

\begin{frame}
  \titlepage
\end{frame}

\begin{frame}{Outline}
    \tableofcontents
\end{frame}

\AtBeginSection[]{
    \begin{frame}{Outline}
        \tableofcontents[currentsection]
    \end{frame}
}

\section{Introduction}

\begin{frame}{Introduction}
  \begin{itemize}
    \item Semantic differences arise across groups, even at the same time.
    \item Word meanings vary across communities and platforms.
    \item We study synchronic semantic variation.
    \item Replace time-based analysis with group-conditioned corpora.
    \item \textbf{Goal:} measure and compare variation using embeddings.
  \end{itemize}
\end{frame}

\begin{frame}{Motivation}
  \begin{itemize}
    \item Most prior work focuses on diachronic semantic change, where word meanings are compared across time using embedding-based methods.
    \item However, semantic variation is also shaped by social and cultural structure, where different communities use the same words differently.
    \item Example: \textit{transformer} refers to a movie concept vs.\ a machine learning model depending on the domain.
    \item \textbf{Gap:} Synchronic, group-based variation is underexplored using the same computational framework.
    \item \textbf{Goal:} Apply and compare multiple methods (word2vec, PPMI, SVD + alignment) to rank words by cross-group variation, find stable anchor words. 
    \item \textbf{Key idea:} If rankings agree (e.g., Spearman correlation), detected variation is robust; disagreement reveals method sensitivity.
  \end{itemize}
\end{frame}

\section{Methodology \& Experiments}

\begin{frame}{Methodology}
    \begin{enumerate}
        \item Fix a time frame, collect corpora from multiple social contexts
        \item Train word embeddings (word2vec, PPMI, SVD) separately per group
        \item Align embedding spaces across groups using Procrustes transformation
        \item Compute cosine distance between aligned embeddings for each word
        \item Compute Spearman correlation between cosine distances produced by different methods
        \item Qualitative analysis of top-ranked words
    \end{enumerate}
\end{frame}

\begin{frame}{Datasets}
    \begin{table}[h]
        \centering
        \begin{tabular}{llr}
            \toprule
            Corpus & Domain & Documents \\
            \midrule
            arXiv CS & Academic (CS) & 16,711 \\
            Yelp     & Restaurant reviews & 369,400 \\
            Ciao     & Product reviews & 48,018 \\
            Reddit   & Online communities & 16M \\
            \bottomrule
        \end{tabular}
        \caption{Corpora used for cross-group comparison}
    \end{table}

    \begin{itemize}
        \item Time window fixed to 2010--2011 (constrained by Ciao availability)
        \item 4 corpora \( \Longrightarrow \) 6 cross-group pairs evaluated per method
    \end{itemize}
\end{frame}

\begin{frame}{Embedding Methods}
    Three distributional methods, each trained independently per corpus:

    \vspace{0.3cm}

    \begin{columns}
        \begin{column}{0.3\linewidth}
            \textbf{word2vec (SGNS)}
            \begin{itemize}
                \item Skip-gram with negative sampling
            \end{itemize}
        \end{column}
        \begin{column}{0.3\linewidth}
            \textbf{PPMI}
            \begin{itemize}
                \item Sparse
                \item High-dimensional
            \end{itemize}
        \end{column}
        \begin{column}{0.3\linewidth}
            \textbf{SVD}
            \begin{itemize}
                \item Dense
                \item Low-dimensional
            \end{itemize}
        \end{column}
    \end{columns}

    \vspace{0.3cm}

    \begin{columns}
        \begin{column}{0.6\linewidth}
            \[ \mathbf{M}^{\text{PPMI}}_{i, j} = \max\!\left(\log \frac{\mathbb{P}(w_i,c_j)}{\mathbb{P}(w_i)\,\mathbb{P}(c_j)} - \alpha,\ 0\right) \]
        \end{column}
        \begin{column}{0.3\linewidth}
            \[ \mathbf{M}^{\text{PPMI}} = \mathbf{U} \boldsymbol{\Sigma} \mathbf{V}^{\top} \]
            \[ \mathbf{w}_i^{\text{SVD}} = (\mathbf{U} \boldsymbol{\Sigma}^\gamma)_i \]
        \end{column}
    \end{columns}

    \vspace{0.3cm}

    \textbf{Alignment:} Procrustes transformation aligns embedding spaces across corpora on shared vocabulary.

    \vspace{0.1cm}

    \textbf{Variation score:} \( d(\vec{w}_A, \vec{w}_B) = 1 - \cos(\vec{w}_A, \vec{w}_B) \), higher means more semantic difference.
\end{frame}

\section{Results \& Conclusion}

\begin{frame}{Spearman Correlation Between Methods}
    \begin{table}[h]
        \centering
        \begin{tabular}{lrrr}
            \toprule
            Pair & w2v vs PPMI & w2v vs SVD & PPMI vs SVD \\
            \midrule
            arXiv vs Yelp   & \(-0.07\) & \(-0.02\) & \(0.44\) \\
            arXiv vs Ciao   & \(-0.03\) & \(0.04\)  & \(0.51\) \\
            arXiv vs Reddit & \(\mathbf{0.29}\)  & \(\mathbf{0.50}\)  & \(0.50\) \\
            Yelp vs Ciao    & \(-0.15\) & \(-0.23\) & \(\mathbf{0.61}\) \\
            Yelp vs Reddit  & \(0.11\)  & \(0.08\)  & \(0.60\) \\
            Ciao vs Reddit  & \(0.14\)  & \(-0.02\) & \(0.54\) \\
            \bottomrule
        \end{tabular}
        \caption{Spearman \( \rho \) on common vocabulary per pair}
    \end{table}

    \begin{itemize}
        \item PPMI vs SVD consistently moderate (\( 0.44 \)--\( 0.61 \))
        \item word2vec agrees with count-based only for similar domains (arXiv--Reddit)
        \item \textbf{PPMI limitation:} distances collapse to 1.0 for distant domains
    \end{itemize}
\end{frame}

\begin{frame}{Anchors via ConShift (Iterative Alignment)}
  \begin{enumerate}
    \item \textbf{Goal:} find stable words across two corpora
    \item Use shared vocabulary between corpora
    \item Align embeddings (rotate source to match target)
    \item Compute shift score (cosine difference)
    \item Select low-shift words as anchors
    \item \textbf{Iterate:} re-align using anchors → update anchors
    \item Stop when anchors stabilize
  \end{enumerate}
\end{frame}

\begin{frame}{Stable Anchors vs Example Variation}

\textbf{Stable words (anchors)}
\begin{itemize}
    \item \textit{difficult}
    \item \textit{larger}
    \item \textit{issues}
\end{itemize}

\vspace{0.5em}

\textbf{Example variations}
\begin{itemize}

    \item \textit{excited  (arXiv vs Reddit)} \hfill \small (domain variation)
    \begin{itemize}
        \item arXiv: ``Every Voronoi cell takes four states: resting, excited, refractory...''
        \item Reddit: ``I'm excited to get into the new decade!''
    \end{itemize}

    \item \textit{valued  (arXiv vs Yelp)} \hfill \small (abstraction variation)
    \begin{itemize}
        \item arXiv: ``...continuous vector-valued function''
        \item Yelp: ``Customer service is just not valued here''
    \end{itemize}

\end{itemize}

\end{frame}

\begin{frame}{Conclusion}
    \textbf{Summary}
    \begin{itemize}
        \item SVD and word2vec both detect cross-group semantic variation; PPMI collapses for distant domains
        \item Moderate Spearman agreement between SVD and word2vec for similar domains (arXiv--Reddit: \( \rho \approx 0.50 \)); near-zero for distant domains
        \item Consistent with \citet{hamilton2016diachronic}: SVD and word2vec are complementary, not equivalent; PPMI is the weakest method
        \item Qualitative analysis confirms interpretable variation (e.g., \textit{excited}, \textit{valued})
    \end{itemize}
\end{frame}

\begin{frame}{Conclusion}
    \textbf{Limitations}
    \begin{itemize}
        \item No ground-truth evaluation, results are qualitative only
        \item Methods cannot distinguish \textit{semantic shift} from \textit{polysemy}: corpora are discrete snapshots, so a word with multiple senses used differently across groups is indistinguishable from a word that has genuinely shifted
        \item PPMI unusable for cross-domain comparison
    \end{itemize}

    \vspace{0.2cm}

    \textbf{Future work}
    \begin{itemize}
        \item Apply to more granular social groups (e.g., subreddits, demographic groups)
        \item Use contextualized embeddings (e.g., BERT) to better separate polysemy from shift
    \end{itemize}
\end{frame}

\begin{frame}{References}
    \tiny
    \nocite{*}
    \bibliographystyle{plainnat}
    \bibliography{bib}
\end{frame}

\end{document}